% T27 source: Примеры базовой модели.tex lines 31--257
\detailedsummaryheading{Elementary M0--M1 Examples}

We first consider four self-contained numerical examples. In every case, the
baseline effort per unit of work is set to $X = 4$~hours.

The idealized AI-assisted completion time is calculated as
\[
  T_{ai}^{(0)} = \frac{1}{P} \sum_{i=1}^{N} Z_i k_i X.
\]

The speedup condition in the idealized model can be written as a parallelism
threshold:
\[
  P > \frac{\sum_{i=1}^{N} Z_i k_i}{\sum_{i=1}^{N} Z_i} = \bar{k}_w.
\]
If the configured parallelism $P$ does not exceed the threshold $\bar{k}_w$,
no speedup is achieved. The condition $P > \bar{k}_w$ is necessary in all
variants of the model; it is sufficient only in the idealized model of
perfectly divisible tasks.

For comparison with the more realistic model of indivisible tasks, we also use
the following lower bound on the makespan:
\[
  T_{ai} \ge \max\left\{ \frac{1}{P} \sum_{i=1}^{N} Z_i k_i X,\; \max_{i \in \overline{1,N}} (Z_i k_i X) \right\}.
\]

The calculations assume $C(P) \equiv 1$ and $\gamma(P) \equiv 1$. Examples
1--4 cover levels M0, where work is treated as perfectly divisible, and M1,
where independent indivisible tasks are considered. The illustrative scenario
additionally considers level M2 with a dependency graph and level M4 with the
developer as a sequential resource. At level M4, each task has a
human-attention component $h_i$ in the decomposition $k_i = a_i + h_i$; the
total human-attention time $H = X \sum_i Z_i h_i$ and the speedup ceiling
$1/\bar{h}_w$ are calculated. A delegated task is divided into initial task
specification, an autonomous agent phase, and final review. From task
specification through acceptance, it retains its assigned agent slot,
including while it waits for developer attention. The schedules are checked
to ensure that human-attention intervals do not overlap and blocked slots
are not reused; their sensitivity to $\gamma(P)>1$ is discussed when the
variants are compared.
Level M3, which describes cross-stream integration overhead $C(P)$, is not
used in the calculations.

\detailedexampleheading{Example 1}

\paragraph{Parameters.}
\begin{table}[ht!]
\centering
\caption{Parameters of Example 1}
\label{tab:example-1-params}
\begin{tabular}{lc}
\toprule
Parameter & Value \\
\midrule
$N$ (number of tasks) & 3 \\
$\{Z_i\}$ (task sizes) & $\{2, 3, 1\}$ \\
$\{k_i\}$ (normalized-duration coefficients) & $\{0.5, 0.7, 0.4\}$ \\
$P$ (parallelism) & 2 \\
$X$ (baseline effort) & 4 h \\
\bottomrule
\end{tabular}
\end{table}

\paragraph{Calculations.}
\begin{align*}
  \sum Z_i &= 2 + 3 + 1 = 6, \\
  \sum Z_i k_i &= 2{\cdot}0.5 + 3{\cdot}0.7 + 1{\cdot}0.4 = 3.5, \\
  T_h &= 6 \cdot 4 = 24~\text{h}, \\
  T_{ai}^{(0)} &= \frac{1}{2} \cdot 3.5 \cdot 4 = 7~\text{h}, \\
  S_E &= \frac{24}{7} \approx 3.43.
\end{align*}

For the indivisible-task model, the AI-assisted task weights are
\[
  \{Z_i k_i X\} = \{4.0, 8.4, 1.6\}~\text{h}.
\]
Under an optimal assignment, the first processor receives the $8.4$-hour task,
and the second receives two tasks with a combined duration of
$4.0 + 1.6 = 5.6$ hours. The attainable makespan is therefore $8.4$ hours, and
the speedup is
\[
  S_E^{\text{ach}} = \frac{24}{8.4} \approx 2.86.
\]

\paragraph{Interpretation.}
All three tasks have $k_i<1$, so AI assistance reduces their modeled effort,
while parallelism $P=2$ amplifies this effect. The idealized model yields a
$3.43\times$ speedup. After task indivisibility is taken into account, the
speedup decreases to $2.86\times$ because the makespan is bounded by the
longest task, whose duration is $8.4$ hours.

\detailedexampleheading{Example 2}

\paragraph{Parameters.}
\begin{table}[ht!]
\centering
\caption{Parameters of Example 2}
\label{tab:example-2-params}
\begin{tabular}{lc}
\toprule
Parameter & Value \\
\midrule
$N$ (number of tasks) & 4 \\
$\{Z_i\}$ (task sizes) & $\{5, 2, 4, 3\}$ \\
$\{k_i\}$ (normalized-duration coefficients) & $\{0.8, 0.6, 0.9, 0.5\}$ \\
$P$ (parallelism) & 3 \\
$X$ (baseline effort) & 4 h \\
\bottomrule
\end{tabular}
\end{table}

\paragraph{Calculations.}
\begin{align*}
  \sum Z_i &= 5 + 2 + 4 + 3 = 14, \\
  \sum Z_i k_i &= 5{\cdot}0.8 + 2{\cdot}0.6 + 4{\cdot}0.9 + 3{\cdot}0.5 = 10.3, \\
  T_h &= 14 \cdot 4 = 56~\text{h}, \\
  T_{ai}^{(0)} &= \frac{1}{3} \cdot 10.3 \cdot 4 \approx 13.73~\text{h}, \\
  S_E &= \frac{56}{13.73} \approx 4.08.
\end{align*}

For the indivisible-task model, the AI-assisted task weights are
\[
  \{Z_i k_i X\} = \{16.0, 4.8, 14.4, 6.0\}~\text{h}.
\]
Under an optimal workload assignment, the processors receive $16.0$, $14.4$,
and $6.0 + 4.8 = 10.8$ hours of work, respectively. The attainable makespan is
$16.0$ hours, and the speedup is
\[
  S_E^{\text{ach}} = \frac{56}{16.0} = 3.50.
\]

\paragraph{Interpretation.}
In Example 2, $k_i<1$ for every task and the parallelism factor is $P=3$. The
idealized speedup of $4.08\times$ is the largest among the first three
examples. Once task indivisibility is taken into account, the makespan is
bounded by the $16.0$-hour task, so the attainable speedup decreases to
$3.50\times$.

\detailedexampleheading{Example 3}

\paragraph{Parameters.}
\begin{table}[ht!]
\centering
\caption{Parameters of Example 3}
\label{tab:example-3-params}
\begin{tabular}{lc}
\toprule
Parameter & Value \\
\midrule
$N$ (number of tasks) & 3 \\
$\{Z_i\}$ (task sizes) & $\{8, 5, 6\}$ \\
$\{k_i\}$ (normalized-duration coefficients) & $\{1.1, 0.7, 0.8\}$ \\
$P$ (parallelism) & 2 \\
$X$ (baseline effort) & 4 h \\
\bottomrule
\end{tabular}
\end{table}

\paragraph{Calculations.}
\begin{align*}
  \sum Z_i &= 8 + 5 + 6 = 19, \\
  \sum Z_i k_i &= 8{\cdot}1.1 + 5{\cdot}0.7 + 6{\cdot}0.8 = 17.1, \\
  T_h &= 19 \cdot 4 = 76~\text{h}, \\
  T_{ai}^{(0)} &= \frac{1}{2} \cdot 17.1 \cdot 4 = 34.2~\text{h}, \\
  S_E &= \frac{76}{34.2} \approx 2.22.
\end{align*}

For the indivisible-task model, the AI-assisted task weights are
\[
  \{Z_i k_i X\} = \{35.2, 14.0, 19.2\}~\text{h}.
\]
Under an optimal assignment, one processor receives the $35.2$-hour task and
the other receives two tasks with a combined duration of
$19.2 + 14.0 = 33.2$ hours. The attainable makespan is $35.2$ hours, and the
speedup is
\[
  S_E^{\text{ach}} = \frac{76}{35.2} \approx 2.16.
\]

\paragraph{Interpretation.}
In Example 3, speedup is retained even though AI assistance increases the
modeled effort of one task ($k_i = 1.1$). This same $35.2$-hour task is the
schedule bottleneck. Consequently, the idealized speedup of $2.22\times$ is
only slightly greater than the attainable value of $2.16\times$; both are
lower than in Examples 1 and 2.

\detailedexampleheading{Example 4}

\paragraph{Parameters.}
\begin{table}[ht!]
\centering
\caption{Parameters of Example 4}
\label{tab:example-4-params}
\begin{tabular}{lc}
\toprule
Parameter & Value \\
\midrule
$N$ (number of tasks) & 3 \\
$\{Z_i\}$ (task sizes) & $\{4, 3, 3\}$ \\
$\{k_i\}$ (normalized-duration coefficients) & $\{2.4, 2.1, 2.3\}$ \\
$P$ (parallelism) & 2 \\
$X$ (baseline effort) & 4 h \\
\bottomrule
\end{tabular}
\end{table}

\paragraph{Calculations.}
\begin{align*}
  \sum Z_i &= 4 + 3 + 3 = 10, \\
  \sum Z_i k_i &= 4{\cdot}2.4 + 3{\cdot}2.1 + 3{\cdot}2.3 = 22.8, \\
  \bar{k}_w &= \frac{22.8}{10} = 2.28, \\
  T_h &= 10 \cdot 4 = 40~\text{h}, \\
  T_{ai}^{(0)} &= \frac{1}{2} \cdot 22.8 \cdot 4 = 45.6~\text{h}, \\
  S_E &= \frac{40}{45.6} \approx 0.88.
\end{align*}

The threshold condition for speedup is not satisfied:
\[
  P = 2 < \bar{k}_w = 2.28.
\]
Consequently, the configured parallelism is insufficient to offset the
increase in modeled task effort under AI-assisted execution.

For the indivisible-task model, the AI-assisted task weights are
\[
  \{Z_i k_i X\} = \{38.4, 25.2, 27.6\}~\text{h}.
\]
For $P=2$, the optimal assignment is $\{38.4\}$ and
$\{25.2 + 27.6 = 52.8\}$ hours. The attainable makespan is therefore $52.8$
hours, and the speedup is lower than the idealized estimate:
\[
  S_E^{\text{ach}} = \frac{40}{52.8} \approx 0.76 < 1.
\]

The makespan lower bound is $\max\{45.6,\; 38.4\} = 45.6$ hours, but it is not
attained: the optimum is $52.8$ hours. This is the only one of the four
examples in which $\max\{W/P,\; \max_i Z_i k_i X\}$ is strictly less than the
optimum. The expression is thus a lower bound rather than a result guaranteed
to be attainable. The gap remains within Graham's guarantee:
$52.8 \le (2 - 1/2) \cdot 45.6 = 68.4$ hours.

The sufficient speedup criterion is not satisfied either. The longest-task
share is $\mu = M_N/T_h = 38.4/40 = 0.96$, and the threshold is
$P - (P-1)\,\mu = 2 - 0.96 = 1.04$, whereas
$\bar{k}_w = 2.28 > 1.04$. Moreover, even the necessary condition
$\bar{k}_w < P = 2$ is violated, consistently with the absence of speedup.

\paragraph{Interpretation.}
In Example 4, all coefficients $k_i$ are substantially greater than one, so
AI assistance markedly increases the modeled effort of every task. With
$\bar{k}_w = 2.28$, the configured parallelism $P=2$ is insufficient to offset
this slowdown. Even the idealized model shows an increase in completion time
from $40$ to $45.6$ hours; after task indivisibility and load imbalance are
taken into account, the makespan increases to $52.8$ hours.

\detailedsummaryheading{Summary Comparison}

\begin{table}[ht!]
\centering
\caption{Summary of the Calculation Results}
\label{tab:summary-results}
\begin{tabular}{lcccc}
\toprule
Metric & Ex.~1 & Ex.~2 & Ex.~3 & Ex.~4 \\
\midrule
$T_h$ (without AI), h & 24.0 & 56.0 & 76.0 & 40.0 \\
$T_{ai}^{(0)}$ (with AI), h & 7.0 & 13.73 & 34.2 & 45.6 \\
Idealized speedup $S_E$ & $3.43\times$ & $4.08\times$ & $2.22\times$ & $0.88\times$ \\
Attainable makespan, h & 8.4 & 16.0 & 35.2 & 52.8 \\
Attainable speedup $S_E^{\text{ach}}$ & $2.86\times$ & $3.50\times$ & $2.16\times$ & $0.76\times$ \\
Speedup condition met? & Yes & Yes & Yes & No \\
\bottomrule
\end{tabular}
\end{table}

\begin{remark}[Final Interpretation]
\begin{enumerate}
  \item Examples 1--3 achieve a speedup both under the idealized formula and
  under the more realistic assignment of indivisible tasks.
  \item The difference between the idealized and attainable speedups is due to
  the longest task: increasing parallelism cannot reduce the makespan below
  its duration.
  \item Example 2 yields the greatest speedup because it combines $k_i < 1$
  for every task with greater parallelism, $P=3$.
  \item Example 3 exhibits a more moderate effect because one of the largest
  tasks has $k_i > 1$ and becomes the critical bottleneck in the schedule.
  \item Example 4 illustrates an adverse scenario: even with $P=2$, the
  parallelism threshold is not satisfied ($P < \bar{k}_w$), so
  $T_{ai}^{(0)} > T_h$ and no speedup is achieved.
\end{enumerate}
\end{remark}
